\documentclass[../main]{subfiles}
\begin{document}
\section{Cotas y Acotación}

\img{images/08/logo}{Acotación en dibujo técnico.}

La acotación es la representación de las dimensiones y otras características de un objeto en el dibujo técnico. Además de las dimensiones, la acotación también representa información adicional (distancias, materiales, referencias, etc.) mediante el uso de líneas, símbolos, figuras y notas. Podemos especificar que la acotación estudia las distintas posibilidades de colocar las medidas en los dibujos de las piezas.

\subsection{Componentes de la cota}

Las cotas se componen de los siguientes elementos:
\img{images/08/ejemplo}{Componentes de cotas: (1) marca inicial (2) línea de cota (3) cifra de cota (4) línea auxiliar de referencia (5) marca final.}
\begin{enumerate}
\item Cifra de cota: es el número que indica la magnitud medida.
\item Línea de cota: es la línea paralela a la arista del objeto que se mide. Sus extremos se indican mediante dos marcas o símbolos.
\item Líneas auxiliares o de proyección de cotas: estas líneas son perpendiculares a lo que están midiendo, y son las que delimitan las líneas de cota.
\item Línea de extensión o de referencia: es una línea auxiliar que va de los extremos de una arista o superficie a los extremos de una línea de cota. Estas líneas sirven para indicar los detalles de la pieza dibujada, y tienen un amplio abanico de posibilidades; pueden acabar en flecha si terminan en el contorno, si lo hacen en círculo si acaban en el interior de la pieza o también pueden acabar en línea cuando terminan en otras líneas de la acotación.
\item Líneas de notas: son aquellas que indican valores o notas adicionales.
Símbolos: son indicaciones gráficas adicionales.
\item Extremidades: hablamos de extremidades cuando nos referimos a los finales de las líneas de cota. Hay una gran variedad de estas, pero la más utilizada es un triángulo isósceles que tiene como ángulo desigual uno de 15°.
     
\end{enumerate}


\newpage
\actividad{Leer, responder y dibujar}
\begin{enumerate}
    \item ¿Cómo se realiza la acotación?
    \item ¿Qué partes componen?
    \item ¿Para qué es necesario el uso de cotas en dibujo técnico?
    \img{images/08/4}{Ejemplo de figura acotada}
    \item Acotar las siguientes figuras:
    \imgfull{images/08/1}{}
    \imgfull{images/08/2}{}
    \img{images/08/3}{}
    \imgfull{images/08/5}{}
\end{enumerate}

\newpage

\newpage

\subsection{¿Qué es acotar?}

En la tarea hay unos dibujos, hay que colocarles las medidas primero medir con la regla.

Se trazan dos lineas, y una linea grande, con flechitas, y se coloca el numero que es la medida, puede ser 2cm, 3cm según lo que mida en el dibujo.

Por ejemplo en ese dibujo hay un 45

Observa esta imagen esta imagen que es un rectángulo es mas claro.

 Se trazan dos lineas auxiliares, y se colocan las flechas y arriba el numero

 la medida del rectangulo es 50


 claro poner las medidas, las flechitas, y las lineas auxiliares, fijate que la cota está compuesta por 3 lineas y 2 flechas.

 Una desde donde comienza a medir, otra donde finaliza, y estan en los extremos las flechas, y en el medio el numero, que indica cuanto mide.

Estas son cotas, y indican medidas, en dibujo técnico.
\img{images/08/cotas}{Cotas.}

habiendo una figura, se mide la figura, que puede ser un dibujo de un objeto real, y se coloca la medida indicando con una cota, siempre usar la regla midiendo desde el cero.

 si esta es tu regla, medis desde el cero y cada rayita es un mm, y cada numero un cm. En dibujo tecnico se usan generalmente los mm o los cm ( centimetros)

 creo que en esta imagen está más claro.
 la actividad lo que hay que hacer es ponerle medidas a los dibujos.
acotarlos
 si lo puede imprimir no es necesario, lo hace en la misma hoja impresa.


\end{document}